\appendix
\section{AHP Pairwise-Comparison Matrices and Their Consistency}
\label{app:ahp}

Several weight vectors in this framework are derived by the Analytic Hierarchy Process. We state the two that back reported quantities, and we report a property of the matrix family that bears on how much the accompanying consistency ratios are worth.

\paragraph{A caveat on consistency ratios} Saaty's consistency ratio detects \emph{in}consistent judgement; it does not detect a matrix constructed from its own answer. If a priority vector $w$ is chosen first and entry $(i,j)$ is then filled in as $w_i / w_j$, the resulting matrix is perfectly consistent, $CR \approx 0$, and the statistic certifies nothing. Such a matrix is rank-one: every row is a scalar multiple of every other. Writing $\sigma_1 \ge \sigma_2$ for its two largest singular values, the ratio $\sigma_2 / \sigma_1$ separates the two cases --- it is near zero for a back-filled matrix and appreciably positive for one carrying genuine judgement disagreement.

\begin{table}[htbp]
\centering
\small
\caption{Consistency diagnostics for the framework's five AHP matrices. $CR$ alone would suggest all five are sound; $\sigma_2/\sigma_1$ shows that three encode a declared priority vector rather than independent pairwise elicitation. The Availability matrix returns a slightly negative $CR$, which is not attainable for a genuinely elicited matrix ($\lambda_{\max} \ge n$) and arises here as floating-point noise around exact consistency.}
\label{tab:ahp-diag}
\begin{tabular}{lccc}
\toprule
\textbf{Matrix} & \textbf{$n$} & \textbf{$CR$} & \textbf{$\sigma_2/\sigma_1$} \\
\midrule
Topic QoS (Eq.~\ref{eq:3}) & 3 & $+0.0158$ & $0.072$ \\
Fault Tolerance & 3 & $+0.0029$ & $0.027$ \\
\midrule
Composite impact ($I_{\text{comp}}$) & 4 & $+0.0011$ & $0.014$ \\
Maintainability & 5 & $+0.0005$ & $0.0006$ \\
Availability & 5 & $-0.0029$ & $0.0006$ \\
\bottomrule
\end{tabular}
\end{table}

\noindent Only the Topic QoS matrix carries a consistency ratio that means what a consistency ratio is supposed to mean, and it is the one matrix a continuous-integration test guards on both sides: \texttt{tests/test\_ahp\_shrinkage.py} asserts $0.005 < CR < 0.10$, the lower bound existing precisely as a non-degeneracy check. We report the remaining vectors as declared constants with documented structure rather than as elicited judgement, and note that this reading costs the framework little: the global sensitivity analysis (Supplementary Section~S1) finds eight of the ten weight constants to have $\mu^* \le 0.023$, so the ranking results in Section~\ref{sec:7} do not rest on these values being optimal.

\begin{table}[htbp]
\centering
\small
\caption{Saaty matrix over the three Topic QoS dimensions backing $\text{QoS}(t)$ in Equation~(\ref{eq:3}). Geometric-mean priority vector $(0.2385, 0.6250, 0.1365)$; $\lambda_{\max} = 3.018$, $CI = 0.0092$, $RI = 0.58$, $CR = 0.0158$. The shipped constants $(0.24, 0.62, 0.14)$ round this within $0.005$.}
\label{tab:ahp-qos}
\begin{tabular}{lccc}
\toprule
 & \textbf{Reliability} & \textbf{Durability} & \textbf{Priority} \\
\midrule
\textbf{Reliability} & $1$ & $1/3$ & $2$ \\
\textbf{Durability} & $3$ & $1$ & $4$ \\
\textbf{Priority} & $1/2$ & $1/4$ & $1$ \\
\bottomrule
\end{tabular}
\end{table}

\begin{table}[htbp]
\centering
\small
\caption{Matrix over the four composite-impact criteria --- reachability loss (RL), fragmentation (FR), throughput loss (TL), flow disruption (FD) --- backing $I_{\text{comp}}(v)$ (Section~\ref{sec:4.3}). Raw priority vector $(0.389, 0.255, 0.255, 0.100)$; after the framework's $\lambda = 0.7$ shrinkage toward a uniform prior, $(0.347, 0.254, 0.254, 0.145)$, which the shipped $(0.35, 0.25, 0.25, 0.15)$ round. Per Table~\ref{tab:ahp-diag} this matrix is rank-one, so it documents the provenance of those constants without independently justifying them.}
\label{tab:ahp-impact}
\begin{tabular}{lcccc}
\toprule
 & \textbf{RL} & \textbf{FR} & \textbf{TL} & \textbf{FD} \\
\midrule
\textbf{RL} & $1$ & $3/2$ & $3/2$ & $4$ \\
\textbf{FR} & $2/3$ & $1$ & $1$ & $5/2$ \\
\textbf{TL} & $2/3$ & $1$ & $1$ & $5/2$ \\
\textbf{FD} & $1/4$ & $2/5$ & $2/5$ & $1$ \\
\bottomrule
\end{tabular}
\end{table}
